\providecommand{\toplevelprefix}{../..}  %
\documentclass[../../book-main_zh.tex]{subfiles}

\begin{document}

\chapter{智能的未知问题与未来方向}
\label{ch:future}

\begin{quote}
    ``{\em 这项研究基于这样一个猜想：学习的每一个方面或智能的任何其他特征，原则上都可以被精确地描述，从而可以制造一台机器来模拟它。
    我们将尝试寻找如何让机器使用语言、形成抽象和概念、解决目前留给人类的各种问题，并自我改进。}''

    $~$\hfill —— 达特茅斯人工智能纲领提案，1956年
\end{quote}
\vspace{5mm}

本书系统地介绍了如何从观测到的高维数据中发展出记忆或经验知识的数学原理和计算机制。
在看似随机的世界中寻求简约规律的能力，是任何智能（无论是自然智能还是人工智能）的基本特征。
我们相信，本书提出的原理和机制是统一且普适的，既适用于动物，也适用于机器。

我们希望本书能帮助读者建立对人工深度神经网络在从高维数据中学习低维分布（的表征）时的功能和作用的严谨理解，
从而彻底澄清围绕现代深度神经网络实践的神秘感。有了这种理解，现有AI模型和系统的能力与局限性便清晰可见：
\begin{enumerate}
    \item 现有的开环模型和系统还远非能够自我学习和自我改进的完整记忆系统。
    \item 现有表征学习的实现仍然是原始且暴力的，在网络架构或优化策略方面远未达到最优。
    \item 现有的AI模型仅仅学习数据分布并进行归纳性的贝叶斯推断，这与高级人类智能有所不同。
\end{enumerate}

本书的目标之一是帮助读者对当前的机器智能技术建立客观、系统的认识，并认清在进一步推动机器智能发展的道路上还有哪些开放问题和挑战。
在这最后一章中，我们将提供一些关于智能的未知领域和未来发展方向的展望。

\section{迈向自主智能：闭环学习？}
从过去十年人工（机器）智能的实践中可以清楚地看到，在给定充足的数据和计算资源的情况下，
人们可以构建一个足够大的模型并对其进行预训练，以学习所有数据的\textit{先验}分布，例如 \(p(\vx)\)。
理论上，这样一个大模型可以记住几乎所有关于世界的现有知识，这些知识被编码在二维图像、三维形状、动态运动以及自然语言等数据中。
正如我们在本书开头所讨论的，这样的大模型扮演着类似于DNA的角色，生命正是利用DNA来记录和传递关于世界的知识。

以这种方式学习到的模型和分布随后可用于生成从同一分布中抽取的新数据样本。人们还可以利用该模型，
在各种条件下利用记住的知识进行推断（例如补全、估计和预测），比如在新的观测 \(\vy\) 下对\textit{后验}分布 \(p(\vx\mid \vy)\) 进行采样。
严格来说，这种推断是归纳性的或统计性的。

任何预训练模型，无论多大，都不能保证它迄今为止学到的分布是完全正确或完整的。如果我们从当前\textit{先验} \(p_{t}(\vx)\) 
中得到的样本 \(\hat{\vx}_{t}\)，或基于\textit{后验} \(p_{t}(\vx\mid \vy)\) 得到的估计 \(\hat{\vx}_{t}(\vy)\) 
与真实值 \(\vx\) 不一致，我们希望基于误差 \(\ve_{t} = \vx_{t} - \hat{\vx}_{t}\) 修正已学习的分布：
\begin{equation}
    p_{t}(\vx) \rightarrow p_{t+1}(\vx) \quad \text{or}\quad
    p_{t}(\vx\mid \vy) \rightarrow p_{t+1}(\vx\mid \vy)。
\end{equation}
这被称为基于误差反馈的误差修正，是自然界中无处不在的持续学习机制。然而，任何开放式（open-ended）模型都缺乏在学到的分布不正确或不完整时对其自身进行修正或改进的机制。
改进当前的AI模型在很大程度上仍然依赖于人类或者外部系统的参与：通过实验、评估和选择进行监督或强化学习。
我们可以将这一过程称为大模型的“人工选择”，类似生命（DNA）进化中的自然选择。

正如我们在本书前面（特别是\Cref{ch:consistent-representations}）所研究的，闭环系统能将内部表征与对外部世界的感知观测对齐。
它们可以不断改进内部学习到的分布及其表征，以实现对外的一致性或自洽性。所以机器智能的研究紧接着的下一步是开发真正的闭环记忆系统
（如\Cref{fig:open-to-closed}所示），使其能够自主且持续地学习任何的数据分布，并基于误差反馈进行不断改进。

因此，从当前流行的端到端训练的开环模型向持续学习的闭环系统过渡
\begin{equation*}
    \text{\textbf{开环}模型} \; \Longrightarrow \;
    \text{\textbf{闭环}系统}
\end{equation*}
是机器真正能够模拟动物和人的大脑在开放世界中学习和应用知识的关键。我们认为
\begin{quote}
    \begin{center}
        \textit{开环模型只适用于一个封闭的世界，无论其多大；\\ 而闭环系统适用于一个开放的世界，无论其多小。}
    \end{center}
\end{quote}
因此，所谓的“通用人工智能”绝不可能仅仅通过记住世界上所有现有的数据和知识来实现。事实上，世界上没有任何知识，无论多少，
是通用的。\footnote{著名的科学哲学家Karl Popper论证过，任何真正的科学理论都是能够被证伪的。} 然而，具有修正自身已有
的记忆和知识不断适应新的环境的能力才是真正通用的。因此，真的智能，只要正确完整的实现了，本身就是通用的。\footnote{我们认为，
AGI应该代表``artificial genuine intelligence''。} 这正是我们在本书开篇讨论的诺伯特·维纳（Norbert Wiener）
在20世纪40年代提出的控制论（Cybernetics）纲领的精髓。

\begin{figure}[t]
    \centering
    \includegraphics[width=0.9\linewidth]{\toplevelprefix/chapters/chapter9/figs/open-to-closed.png}
    \caption{从开放式深度网络到闭环系统。}
    \label{fig:open-to-closed}
\end{figure}

\section{迈向自然智能：超越反向传播？}
近年来机器智能的实践使许多人相信，必须构建一个单一的大模型来学习所有数据的分布并记住所有知识。
尽管这在技术上是可行的，但这种解决方案可能远非必要或高效。正如我们从训练深度网络中所知，
唯一已知的大规模训练此类网络的可扩展方法是反向传播（BP）\cite{Back-Prop}。
尽管BP提供了一种通过在整个模型中反向传播梯度信号来修正误差的方法，但它仍然相当暴力，
与自然界的学习方式大相径庭：由于成本高昂，BP是自然界无法承受的选项。并且由于生物物理机制的限制，自然界根本无法实现它。

更一般地说，除非我们也理解智能如何被高效地实现，否则我们无法真正理解智能。
也就是说，我们需要解决实现与达成智能目标相关的机制的计算复杂度问题。
从历史上看，我们对（机器）智能的理解经历了几个阶段，从不可计算的柯尔莫哥洛夫（Kolmogorov）复杂度
到香农（Shannon）熵，从图灵（Turing）的可计算性到后来对易解性（tractability）的理解，\footnote{如果一个问题
存在一个复杂度随问题规模呈多项式增长的算法，我们就称该问题是易解的（tractable）。}
再到现代高维数据分析\cite{Wright-Ma-2022}和人工智能实践中对算法可扩展性（scalability）的强烈强调。这种演变可以总结如下：
\begin{equation*}
    \text{\textbf{不可计算}} \;
    \Longrightarrow \; \text{\textbf{可计算}} \;
    \Longrightarrow \; \text{\textbf{易解}} \; \Longrightarrow \;
    \text{\textbf{可扩展}}。
\end{equation*}
在很大程度上，深度学习和反向传播的成功与普及，正是因为它们在现代计算平台（如GPU）上提供了一种可扩展的实现方式，能用于处理和压缩海量数据。
尽管如此，这种实现的成本仍然远远高于自然界实现智能的方式。

机器智能的效率仍有巨大的提升空间，使其能够模拟自然智能的效率，这应该比当前暴力的实现方式高出好几个数量级。
为此，我们需要发现新的学习架构和优化机制，使其能够在自然物理条件和资源约束下学习数据分布，
类似于自然界中智能生物所面临的情况——例如，无需一次性访问所有数据，也无需（通过BP）同时更新所有模型参数。

本书提出的原则性框架和方法可以指导我们发现这些新的架构和机制。这些新架构和机制应能实现实时在线的持续学习，
并应能通过高度局部化和稀疏的前向或后向优化进行更新。\footnote{对于学习分布而言，这些要求的一个简单实例是实现PCA的方法，
即\Cref{ch:consistent-representations}中介绍的在线PCA方法。一旦学到了这样的线性模型，
通过有噪声的观测进行在线预测模型这一任务，就可以通过最小方差这一类方法实现，例如维纳滤波\cite{Wiener-1942,Wiener-1949}。
对于动态的线性系统，这就成为了有名的卡尔曼滤波，见\cite[Appendix B]{MaY2003} 或\cite{peng2025mathematics}。}

\begin{figure}[t]
    \centering
    \includegraphics[width=0.99\linewidth]{\toplevelprefix/chapters/chapter9/figs/loops.png}
    \caption{大脑皮层的猜想架构。大脑皮层是一个大规模并行和分布式的自编码系统，由闭环自编码器的层级结构组成。它们从输入中提取信息，并最大化其输出表征的信息增益。这些过程在层级结构的多个层面上进行，具有不同层次的信息抽象。}
    \label{fig:loops}
\end{figure}
正如我们从神经科学中所了解到的，我们大脑的皮层由数以万计的皮层柱组成\cite{Hawkins-2021}。
所有皮层柱都具有相似的物理结构和功能。它们是高度并行和分布式的，相互之间有着一些很稀疏的连接。
因此，我们认为，为了开发一个更具可扩展性和结构化的记忆系统，我们需要考虑模拟大脑皮层的架构。
\Cref{fig:loops}展示了这样一种假设的架构：一个大规模分布式和层级化的系统，
由许多在很大程度上并行的闭环自编码模块组成。\footnote{这种假设架构的元素以各种形式存在于文献中，
例如预测编码网络\cite{millidge2022predictive}，它具有利用高效局部更新的优化策略。}
这些模块学习对不同的感觉模态或来自每种感觉模态的数据的不同投影进行编码。
我们在\Cref{ch:conditional-inference}的\Cref{sec:measurement-self-consistency}中的讨论表明，
这种对低维分布的并行感知和学习在理论上是可行的。在这之后，基于较低层级的输出表征可以继续学习较高层级的（有损）自编码，
以对较低层级学到的表征发展出更稀疏、更高层级的“抽象”。

\Cref{fig:loops}中所示的分布式、层级化和闭环系统架构具有大脑皮层的许多特征。
这种系统架构可能会比当前的单一的大模型架构开启更多的可能性。它使得探索更高效的学习和优化机制成为可能，
并导致对学到的数据分布和知识形成更结构化、模块化的组织，更接近生物大脑里的记忆系统。
这将使我们能够把机器智能的实现带入进化的下一个阶段：
\begin{equation*}
    \text{\textbf{不可计算}}
    \Longrightarrow \text{\textbf{可计算}}
    \Longrightarrow \text{\textbf{易解}} \Longrightarrow
    \text{\textbf{可扩展}}
    \Longrightarrow
    \text{\textbf{自然}}。
\end{equation*}

\section{迈向科学智能：超越图灵测试？}

\subsection{自然界中智能的进化}
正如我们在本书开篇\Cref{ch:intro}中所讨论的，自然界中的智能经历了多个阶段的进化，并表现为四种截然不同的形式：
\begin{equation*}
    \text{\textbf{物种进化智能}} \;
    \Longrightarrow \; \text{\textbf{个体成长智能}} \;
    \Longrightarrow \; \text{\textbf{社会发展智能}} \;
    \Longrightarrow \; \text{\textbf{科学进步智能}}.
\end{equation*}
所有形式的智能都有一个共同的目标（和特征），即从外部世界感知到的高维数据学习到其中低维的结构，从而转化成有用的知识。
尽管如此，这些不同的智能形式在所采用的具体编码方案、所编码的信息、用于学习和改进的优化计算机制，以及这些机制的物理实现方面有所不同。
使用本书中介绍的概念和术语，这四个阶段的智能在以下三个方面存在差异：
\begin{enumerate}
    \item 用于编码信息或知识的\textit{编码字典}或方案。
    \item 使用码本学习和编码的\textit{信息}或知识。
    \item 用于创建和修正编码信息或知识的\textit{优化机制}。
\end{enumerate}
表 \ref{tab:four-stage-intelligence} 总结了它们的主要特征。
\begin{table}[h]
    \centering
    \resizebox{1\columnwidth}{!}
    {
        \begin{tabular}{| c | c | c | c | c |}
            \hline & \textbf{物种进化} & \textbf{个体成长} &
            \textbf{社会发展} & \textbf{科学进步}\\
            \hline

            \textbf{编码字典}  & DNA & 神经元/大脑 & 自然语言 & 数学抽象 \\ [0.5ex]
            \hline

            \textbf{信息内容} & 基因 & 记忆 & 经验知识 & 科学规律\\ [0.5ex]
            \hline
            \textbf{优化机制} & 强化学习 & 误差反馈 &
            归纳试错 & 假设与证伪 \\  [0.5ex]
            \hline
        \end{tabular}
    }
    \caption{自然界智能四个阶段的主要特征。}
    \label{tab:four-stage-intelligence}
\end{table}

正如我们现在所知，人类在智能方面实现了两次质的飞跃。
\begin{enumerate}
    \item 第一次是口头和书面语言的发展，这使得人类能够跨个体甚至跨代分享和传递学到的知识，就像身体里的DNA一样。
    \item 第二次是大约三千年前数学和形式逻辑的发展，它们成为了现代科学的精确语言和演绎方法。
\end{enumerate}
数学语言使我们摆脱了仅仅以经验形式从观测中总结归纳知识的局限，
并允许我们将知识形式化为可以通过数学演绎或实验来验证或证伪的理论。
通过提出假设、逻辑演绎以及实验验证或证伪，我们现在可以主动发现和发展新知识，
这远远超出了仅从观测数据分布中所能学到的内容。\footnote{例如，因果关系不能仅从分布中学习\cite{Pearl-2009}。}

尽管这些不同阶段的智能看似存在差异，但智能的进化具有一个共同特征：
{\em 它们都使用某种形式的反馈\footnote{对错误的损失和对成功的奖励}来创建和修正学到的信息，
并且在进化过程中，反馈和修正变得越来越频繁和高效——从开环的强化学习（自然选择）到闭环自监督（大脑记忆），从集体试错到主动假设与证伪。}

\subsection{从归纳智能到演绎智能}
正如在引言（\Cref{ch:intro}）中所讨论的，1956年的“人工智能”（AI）纲领正是旨在研究如何赋予机器科学智能，即数学抽象、逻辑推断和解决问题等被认为是区分人类与动物的高级功能：
\begin{equation*}
    \text{\textbf{低级}（动物）智能} \; \Longrightarrow \;
    \text{\textbf{高级}（人类）智能.}
\end{equation*}
正如我们在本书中反复澄清的那样，过去十多年来机器智能的大多数技术进步，虽然是在“AI”的名义下进行的，
但实际上更密切相关的是让机器模仿动物和人类共有的低级智能机制，包括物种进化智能、个体发育智能和社会智能。
在这些层面上，智能创造的是{\em 经验}记忆或知识，并在使用时进行本质上主要是{\em 归纳}的推断。
本书介绍的原理和方法旨在揭示这些层面智能背后的科学目标和计算机制。它们为以下论点提供了强有力的理论依据和计算证据：
\begin{quote}
    \centering
    \textit{机器是可以从数据中学习经验知识并进行归纳推断。}
\end{quote}
然而，到目前为止，还没有严谨的理论或科学证据表明，仅靠这些机制就足以实现最初的AI纲领的真正主旨：理解人类的高级智能的机理并希望机器实现它。

事实上，时至今日，我们对于如何严谨地验证或证明一个系统是否真正具备这种高级智能知之甚少，
尽管图灵测试（也称为模仿游戏）\footnote{在图灵的提议中，评估者或询问者是人类。然而，大多数人类评估者的科学训练和知识有限，
他们的结论可能是主观的。}早在1950年就被提出\cite{Turing-1950}，作为解决以下根本问题的初步尝试：
\begin{quote}
    \centering
    \textit{机器能思考吗？}
\end{quote}
很长一段时间以来，人们认为没有必要对这种测试进行精确定义，因为机器的能力远低于人类（甚至动物）。
然而，鉴于近年来的技术进步，许多模型和系统现在声称达到甚至超越了人类智能。因此，现在是时候制定一个科学且可验证的图灵测试定义了，
即一个系统且严谨的协议，用于评估和证明任何所谓“智能”的模型或系统的智能水平。

正如图灵在其原始论文中所论述的那样，定义什么是“思考”相当困难。这里我们可以从更具体的问题开始，
尝试确定一个系统是否真正“理解”某些概念或知识。例如，我们如何严谨地验证一个智能模型或系统是否真正理解了诸如以下的抽象概念：
\begin{enumerate}
    \item 两个量之间绝对相等的概念，
    \item 数字（自然数、有理数或实数）的概念，
    \item 无穷的概念，
    \item 以及数学归纳法的概念。
\end{enumerate}
还是它仅仅记住了大量关于这些概念的例子？学到这些例子的一些经验规律。请注意，最先进的大语言模型仍然在诸如“3.11比3.9大还是小？”这样简单的数学问题上挣扎。
\footnote{一些模型通过有针对性的工程设计纠正了它们对这类问题的回答，或者结合了额外的验证机制，
在推断过程中验证即时答案并进行纠正。然而，我们留给读者一个练习：
严谨地测试任何最先进的语言模型本身是否真正理解数字（自然数、有理数、实数和复数）的概念及其相关的算术。}

通过学习本书，我们应该知道，现代生成式AI系统主要依赖贝叶斯（因此是通过归纳的方法）推断来生成答案。
机器也能真正掌握数学归纳法或一般的逻辑演绎，并理解它与贝叶斯归纳的区别吗？
我们如何验证一个系统是否真正理解逻辑和数学演绎的规则，并能严谨地应用它们。
或者，它是否仅仅记住了大量此类演绎的实例（例如思维链数据），并学习和利用了它们的分布结构？
因此，我们需要回答下面这个未知的问题：
\begin{quote}
    \centering
    \textit{\textbf{记忆}与\textbf{理解}之间有区别吗？}
\end{quote}
只有明确和了解了这种区别，我们才有可能尝试对“{\em 机器能理解吗？}”这个问题给出一个有意义的答案。

也许更重要的是，如何验证一个自称智能的系统是否具备{\em 创造新知识和修正自身知识}的能力，例如创造新的数学概念、
发现新的物理定律和因果关系？理解与自我创造之间的联系可能是相当根本的，正如物理学家理查德·费曼（Richard Feynman）的一句名言：
\begin{quote}
    \centering
    \textit{“我不能自己创造出来的，我就无法理解。”}
\end{quote}

\subsection{智能的科学测试}
因此，现在是时候建立一套科学合理的评估方法，以确定一个系统的智能能力属于以下哪一类了：
\begin{enumerate}[label={Level-\arabic*.}, leftmargin=*]
    \item 仅仅学习了一些承载知识的数据的分布，并能够重新生成它们；
    \item 能够从新的观测和经验中自主且持续地学习新知识或修正现有知识；
    \item 能够真正理解和创造抽象概念、规则和定律，并知道如何以演绎的方式正确应用它们；
    \item 能够生成新的科学假设或数学猜想，并基于实验或逻辑演绎对其进行验证或证伪。
\end{enumerate}

\begin{figure}[t]
    \centering
    \includegraphics[width=\linewidth]{\toplevelprefix/chapters/chapter9/figs/tests.png}
    \caption{针对不同级别或类型智能能力的三种测试：针对基础智能的维纳测试，针对人类高级别智能的图灵测试，以及针对科学家级别智能的波普尔测试。}
    \label{fig:three-tests}
\end{figure}

为了评估和区分这些不同级别的智能，我们建议可能至少需要三种不同的测试，如\Cref{fig:three-tests}所示：
\begin{enumerate}
    \item \textit{诺伯特·维纳测试（The Norbert Wiener Test）：}确定一个系统是否能够自我修正并自行创造新的（经验）知识，还是仅仅通过外部强化或监督被动地接收和更新信息；
    \item \textit{艾伦·图灵测试（The Alan Turing Test）：}确定一个系统是否能够创造和理解抽象概念与知识，还是仅仅学习和记忆样本的统计特征并将其用于贝叶斯推断；
    \item \textit{卡尔·波普尔测试（The Karl Popper Test）：}确定一个系统是否能够通过基于实验一致性或逻辑自洽性形成并验证假设或理论，从而主动探索和创造新的抽象知识。
\end{enumerate}
我们认为，对于此类测试，评估者或询问者不应是任意的单个人或群体，而应是遵循具有科学性的合理协议和流程，具有科学权威的实体，以便得出的结论严谨且值得信赖。

正如我们在本书中所看到的，\textit{压缩}在学习或创造经验知识中发挥了最本质作用。
它是识别经验数据分布并以更好的表征组织编码其信息的指导原则和普适机制。在很大程度上，
它解释了迄今为止“人工智能”的绝大部分实践。因此这些实践实现的主要是归纳性的智能。
未来研究的一个重点问题是，\textit{仅仅依靠压缩}是否足以实现上述所有更高级别的具有演绎能力的智能。
特别是，数学抽象、因果推断、假设生成以及逻辑推理和演绎，是压缩的自然延伸还是具有本质的不同？
通过压缩识别经验数据分布与通过抽象形成高级概念和理论之间，是否存在某种根本区别？因此，
我们需要回答下面这个未知的问题：
\begin{quote}
    \centering
    \textit{\textbf{压缩}与\textbf{抽象}之间有区别吗？}
\end{quote}
只有了解了这种区别，我们才有可能尝试对“{\em 机器能抽象或创造吗？}”这个问题给出一个有意义的答案。

在很大程度上，科学——及其相关的码本，数学——可以被视为最高级的智能形式，是受过教育和启蒙的人类所独有的，类似于一种高级艺术形式。哲学家卡尔·波普尔（Karl Popper）爵士曾提出：
\begin{quote}
    \centering
    \textit{``科学可以被描述为一门对知识进行系统性的极度简化的艺术。''}
\end{quote}
我们相信，揭示并理解这种更高层次智能背后的数学原理和计算机制，并通过机器去成功复现，将是科学、数学和计算共同的最终前沿！

\end{document}